\chapter{Liver Resection (Hepatectomy)}
\section*{Overview}
Liver resection removes a diseased portion of the liver while preserving enough functioning liver tissue.
\section*{Common Instruments and Equipment}
The exact setup varies by technique and setting. Common items may include:
\begin{itemize}[nosep,leftmargin=*]
  \item open or minimally invasive liver surgery instrument set
  \item ultrasonic dissector or energy device for parenchymal transection
  \item intraoperative ultrasound and vascular control instruments
\end{itemize}

\section*{Why It Is Done}
It may treat selected primary liver cancers, metastatic tumors, benign tumors, or localized injury and disease.
\section*{How It Is Performed}
The operation may be open, laparoscopic, or robotic. The surgeon controls liver vessels and bile ducts, removes the planned segment or lobe, and controls the cut liver surface.
\section*{Contraindications}
No single contraindication applies to every patient; the treating team assesses whether the procedure is suitable for the individual and the reason for treatment.

\section*{Precautions}
The care team reviews relevant health conditions, medicines, allergies, and procedure-specific risks, and explains any changes needed for the individual.

\section*{Risks and Recovery}
Bleeding, bile leak, infection, liver failure, blood clots, and recurrence are possible. Liver function, imaging, pathology, and oncology follow-up are required.
\section*{References}
\begin{enumerate}\item National Cancer Institute. Liver Cancer Surgery. 2025.\item Current Medical Diagnosis and Treatment. McGraw-Hill; 2025.\end{enumerate}
\clearpage
